%% ma: L11-B17-0060
%% lop: 11
%% chuong: 5
%% bai: 17
%% dang: 11.17.3
%% loai: TN4
%% muc: VD
%% dap_an: "C"
%% nguon: "tai-lieu-hoc-tap-mon-toan-lop-11-theo-chuong-trinh-giao-duc-pho-thong-2018.pdf (Lớp Toán Thầy Duy, Phạm Lê Duy, THPT Chu Văn An), Chương 3 Bài 3, Phần 1, Câu 40 (tr.411)"
%% kien_thuc: "Lớp 11 Bài 17 (hàm số liên tục)"
%% trang_thai: chinh-thuc
%% ghi_chu: "Đã duyệt 03/11/2026."
\begin{ex}
Cho số thực $a,b,c$ thỏa mãn $\begin{cases}-8+4a-2b+c>0\\8+4a+2b+c<0\end{cases}$. Số giao điểm của đồ thị hàm số $y=x^3+ax^2+bx+c$ và trục $Ox$ là
\choice{$2$.}{$0$.}{\True $3$.}{$1$.}
\loigiai{\sloppy Đặt $f(x)=x^3+ax^2+bx+c$, liên tục trên $\mathbb R$. Giả thiết cho $f(-2)>0$ và $f(2)<0$. Vì $\lim\limits_{x\to-\infty}f(x)=-\infty$ nên có $u<-2$ với $f(u)<0$, suy ra $f(u)f(-2)<0$; $f(-2)f(2)<0$; vì $\lim\limits_{x\to+\infty}f(x)=+\infty$ nên có $v>2$ với $f(v)>0$, suy ra $f(2)f(v)<0$. Phương trình có ít nhất $3$ nghiệm trong ba khoảng rời nhau $(u;-2)$, $(-2;2)$, $(2;v)$, mà đa thức bậc ba có tối đa $3$ nghiệm. Vậy có đúng $3$ giao điểm.}
\end{ex}
